\documentclass[10pt]{article}
\usepackage[T1]{fontenc}\usepackage[utf8]{inputenc}\usepackage{textcomp}
\usepackage[scaled]{helvet}\renewcommand{\familydefault}{\sfdefault}
\DeclareUnicodeCharacter{03B5}{\ensuremath{\varepsilon}}\DeclareUnicodeCharacter{03B2}{\ensuremath{\beta}}
\DeclareUnicodeCharacter{03B7}{\ensuremath{\eta}}\DeclareUnicodeCharacter{03B1}{\ensuremath{\alpha}}
\DeclareUnicodeCharacter{03C7}{\ensuremath{\chi}}\DeclareUnicodeCharacter{2212}{\ensuremath{-}}
\DeclareUnicodeCharacter{2265}{\ensuremath{\geq}}\DeclareUnicodeCharacter{2264}{\ensuremath{\leq}}
\DeclareUnicodeCharacter{2192}{\ensuremath{\rightarrow}}
\usepackage[a4paper,margin=10mm]{geometry}
\usepackage{graphicx}
\pagestyle{empty}\setlength{\parindent}{0pt}\setlength{\parskip}{3pt}
\begin{document}
{\small\textbf{Fig. S5}\hfill }\par\vspace{2mm}
\begin{center}\includegraphics[width=174mm,height=200mm,keepaspectratio]{/mnt/data/shahar/gwas_barley/morexV3_analysis/barley_for_publication_project/08_USED_creating_figures/Figure_ESM_site_BLUPs/ESM_site_BLUPs.pdf}\end{center}
\vspace{1mm}{\small \textbf{Fig. S5} Genotype BLUPs of the four nutritional traits (centered; \% of grain) of the 290 accessions across the 29 sampling sites, ordered and colored by ecological region (HZ1, HZ2 and HZ3, the North–Coast, North–Desert and Coast–Desert intermediate zones). Boxes, median and interquartile range; whiskers, 1.5 × the interquartile range; points, accessions (ten per site). Sites are numbered as in Table S1 (every fifth site labeled; site 4 was not analyzed)\par}
\end{document}
